\subsection{What We Can Defend}

The most defensible claim at this stage is qualitative rather than exact:
\begin{enumerate}[nosep]
  \item barriers tend to rise when current-state reinforcement is stronger;
  \item barriers tend to rise when the target state is ambiguous;
  \item barriers tend to fall when relevant preparation has already been done;
  \item barriers depend on current conditions \emph{and} transition history.
\end{enumerate}

\begin{definition}[Preparation and ambiguity]
Let $U(s_1)$ denote \textbf{target ambiguity}, meaning the unresolved
specification burden of entering $s_1$. Let $P$ denote \textbf{preparation},
meaning the extent to which entry-relevant work has already been performed
before the transition begins.
\end{definition}

\begin{remark}
\textbf{Monotone barrier scaffold.}
For modelling purposes, the transition difficulty may be treated as depending
on a function
\[
  \DV(s_0 \to s_1) \approx g\!\bigl(\widehat{\FI}(s_0), U(s_1), P, H\bigr),
\]
where $H$ summarizes transition-relevant history, and where $g$ is assumed to
be increasing in current-state resistance and ambiguity, and decreasing in
preparation.
\end{remark}

\begin{discussion}
This assumption is deliberately weaker and more defensible than the chapter's
earlier barrier-decomposition formula. It keeps the business logic that later
chapters need: if the present state is strongly reinforced, if the target state
is vague, and if preparation is weak, the transition becomes harder. But it no
longer pretends that the manuscript has identified a specific additive law with
known coefficients.
\end{discussion}

\begin{remark}
The symbol ``$\approx$'' should be read conservatively here. It does not mean
statistical fit or asymptotic approximation. It means only that the barrier
construct is being represented by a modelling dependence on those named inputs.
\end{remark}
